\section{Conclusions and Future Work}
\label{sec:conclusions}

This article presents a structured integration of the Data Privacy Vocabulary with the Open Digital Rights Language to enable the representation of legally grounded, machine-readable usage control policies in data spaces. By defining a comprehensive mapping between DPV taxonomies and ODRL constructs, and formalising it as an interoperable ODRL profile, the work bridges the gap between legal semantics and enforceable policy representations. The resulting approach allows policies to capture permissions, prohibitions, and obligations enriched with domain-specific concepts such as purposes, legal bases, risks, and technical measures, thereby supporting the expression of compliance requirements derived from data protection and governance frameworks. Overall, this contribution provides a reusable and extensible foundation for aligning legal, organisational, and technical perspectives in the design of trustworthy and interoperable data sharing ecosystems.

Future work will focus on the implementation and evaluation of this profile within interoperable ODRL policy engines~\cite{slabbinck_interoperable_2025} to support automated policy interpretation, validation, and enforcement across heterogeneous data space infrastructures. This includes extending existing ODRL engines to natively support DPV-based semantics and integrating SHACL-based validation mechanisms for constraint checking. Additionally, further research is needed to assess scalability and interoperability across different data space implementations, as well as to explore integration with data negotiation protocols and trust frameworks. These efforts will contribute to operationalising the proposed profile and advancing practical, standards-based solutions for legally compliant usage control in real-world data spaces.